% Options for packages loaded elsewhere
\PassOptionsToPackage{unicode}{hyperref}
\PassOptionsToPackage{hyphens}{url}
\PassOptionsToPackage{dvipsnames,svgnames,x11names}{xcolor}
\documentclass[
]{article}
\usepackage{xcolor}
\usepackage[margin = 0.5in]{geometry}
\usepackage{amsmath,amssymb}
\setcounter{secnumdepth}{-\maxdimen} % remove section numbering
\usepackage{iftex}
\ifPDFTeX
  \usepackage[T1]{fontenc}
  \usepackage[utf8]{inputenc}
  \usepackage{textcomp} % provide euro and other symbols
\else % if luatex or xetex
  \usepackage{unicode-math} % this also loads fontspec
  \defaultfontfeatures{Scale=MatchLowercase}
  \defaultfontfeatures[\rmfamily]{Ligatures=TeX,Scale=1}
\fi
\usepackage{lmodern}
\ifPDFTeX\else
  % xetex/luatex font selection
\fi
% Use upquote if available, for straight quotes in verbatim environments
\IfFileExists{upquote.sty}{\usepackage{upquote}}{}
\IfFileExists{microtype.sty}{% use microtype if available
  \usepackage[]{microtype}
  \UseMicrotypeSet[protrusion]{basicmath} % disable protrusion for tt fonts
}{}
\makeatletter
\@ifundefined{KOMAClassName}{% if non-KOMA class
  \IfFileExists{parskip.sty}{%
    \usepackage{parskip}
  }{% else
    \setlength{\parindent}{0pt}
    \setlength{\parskip}{6pt plus 2pt minus 1pt}}
}{% if KOMA class
  \KOMAoptions{parskip=half}}
\makeatother
\usepackage{graphicx}
\makeatletter
\newsavebox\pandoc@box
\newcommand*\pandocbounded[1]{% scales image to fit in text height/width
  \sbox\pandoc@box{#1}%
  \Gscale@div\@tempa{\textheight}{\dimexpr\ht\pandoc@box+\dp\pandoc@box\relax}%
  \Gscale@div\@tempb{\linewidth}{\wd\pandoc@box}%
  \ifdim\@tempb\p@<\@tempa\p@\let\@tempa\@tempb\fi% select the smaller of both
  \ifdim\@tempa\p@<\p@\scalebox{\@tempa}{\usebox\pandoc@box}%
  \else\usebox{\pandoc@box}%
  \fi%
}
% Set default figure placement to htbp
\def\fps@figure{htbp}
\makeatother
\setlength{\emergencystretch}{3em} % prevent overfull lines
\providecommand{\tightlist}{%
  \setlength{\itemsep}{0pt}\setlength{\parskip}{0pt}}
\usepackage{color}
\usepackage{framed}
\setlength{\fboxsep}{.8em}

\usepackage{tcolorbox}

\tcbset{
  mystyle/.style={
    before skip=-10pt, % Adjust this value
    after skip=-10pt   % Adjust this value
  }
}
\newtcolorbox{mynoteblock}[1][]{mystyle, #1}

\newtcolorbox{blackbox}{
  colback=black,
  colframe=orange,
  coltext=white,
  boxsep=5pt,
  arc=4pt}
  
\newtcolorbox{ltbluebox}{
  colback=blue!5!white,
  colframe=blue!75!black,
  boxsep=5pt,
  arc=4pt}
  
\newtcolorbox{ltgreenbox}{
  colback=green!5!white,
  colframe=blue!75!black,
  boxsep=5pt,
  arc=4pt}
  
\newtcolorbox{ltpurplebox}{
  colback=purple!5!white,
  colframe=blue!75!black,
  boxsep=5pt,
  arc=4pt}
  


\newenvironment{cols}[1][]{}{}

\newenvironment{col}[1]{\begin{minipage}[t]{#1}\ignorespaces\vspace{0pt}}{%
\end{minipage}
\ifhmode\unskip\fi
\aftergroup\useignorespacesandallpars}

\def\useignorespacesandallpars#1\ignorespaces\fi{%
#1\fi\ignorespacesandallpars}

\makeatletter
\def\ignorespacesandallpars{%
  \@ifnextchar\par
    {\expandafter\ignorespacesandallpars\@gobble}%
    {}%
}
\makeatother

\usepackage{awesomebox}
\setlength{\aweboxvskip}{1mm}


\pagenumbering{gobble}

\usepackage{enumitem}
\setlist{itemsep=1pt, parsep=0pt} %# Adjust vertical spacing
\setlist[itemize,1]{leftmargin=*, labelsep=0.3em} %# Adjust horizontal spacing

\usepackage{bookmark}
\IfFileExists{xurl.sty}{\usepackage{xurl}}{} % add URL line breaks if available
\urlstyle{same}
\hypersetup{
  colorlinks=true,
  linkcolor={Maroon},
  filecolor={Maroon},
  citecolor={Blue},
  urlcolor={blue},
  pdfcreator={LaTeX via pandoc}}

\author{}
\date{\vspace{-2.5em}}

\begin{document}

\vspace{-1cm}

\section{MAFMC Indicators Project: Workshop 1
Results}\label{mafmc-indicators-project-workshop-1-results}

Workshop objective: Identify and prioritize Council management decisions
that could be supported or improved by operational use of ecosystem
indicators.

\subsection{December 10, 2025 Workshop Attendance: 40
total}\label{december-10-2025-workshop-attendance-40-total}

\begin{noteblock}
\emph{In person:} 13 participants from the Council, Council Advisory
Panels, Scientific and Statistical Committee, Council staff, and
Atlantic States Marine Fisheries Commission staff, plus 4 from the
project team.

\emph{Online:} 21 participants from the Council, Council Advisory
Panels, Council staff, Northeast Fisheries Science Center staff, and
Project Oversight Team, plus 2 from the project team.

\end{noteblock}

\begin{ltgreenbox}

\subsection{Overview of Workshop Discussions and
Outcomes}\label{overview-of-workshop-discussions-and-outcomes}

\textbf{Goal:} Which management decisions should this project focus on?
Why?

\textbf{Discussion 1:} Reviewed and updated the list of management
decisions

\textbf{Discussion 2:} Identified 4 priority decisions

\textbf{Discussion 3:} Provided initial feedback on SMART indicator
review for further development

\end{ltgreenbox}

\begin{cols}

\begin{col}{0.60\textwidth}

\subsection{Discussion 1: Clarify
Decisions}\label{discussion-1-clarify-decisions}

\begin{itemize}
\tightlist
\item
  Frame management decisions to reflect both restriction and
  liberalization
\item
  Consider distinctions between:

  \begin{itemize}
  \tightlist
  \item
    Pre-season and inseason decisions
  \item
    Short term and longer term planning horizons
  \item
    Routine and emergency contexts
  \end{itemize}
\item
  Consider tradeoffs between objectives in decisions
\item
  Consider uncertainty: indicators should reduce it
\item
  Consider stages in the process where information is most useful
\item
  Additions: Quota management, Decisions regarding emerging fisheries,
  MMPA and ESA objective interactions with EFH, Strategic planning and
  Annual Action Plan, Shared governance
\end{itemize}

\begin{ltgreenbox}

\subsection{Discussion 2: Priority
Decisions}\label{discussion-2-priority-decisions}

\begin{itemize}
\tightlist
\item
  \textbf{Risk policy:} consider risk indicators in addition to stock
  status
\item
  \textbf{Catch/quota setting:} communicate information available/used
  throughout the process. Consider Assessment, SSC ABC, Monitoring
  Committee, and Council processes
\item
  \textbf{Setting measures:} consider multispecies and multifleet
  context in the specifications process. Objective: maximize
  opportunities across multiple management bodies.
\item
  \textbf{Spatial management measures:} consider changing distributions
  of fish and fleets within the framework process
\end{itemize}

\paragraph{For all selected
decisions:}\label{for-all-selected-decisions}

\begin{itemize}
\tightlist
\item
  Evaluate performance of current management against stated objectives
  to aid indicator selection
\item
  Evaluate impact of data availability or degradation on uncertainty and
  indicator performance
\end{itemize}

\end{ltgreenbox}

\subsection{Discussion 3: Indicator
Tools}\label{discussion-3-indicator-tools}

\begin{itemize}
\tightlist
\item
  SMART Indicator framework and tools a promising approach
\item
  Add information on state level surveys and other information used in
  Habitat Assessment
\item
  Design indicator explorer tools to bring information into many
  different Council products and processes\\
\end{itemize}

\end{col}

\begin{col}{0.02\textwidth}
~

\end{col}

\begin{col}{0.38\textwidth}

\emph{Project Steps: This workshop focused on Management Objectives and
Decisions, and outlined additional steps}

\includegraphics[width=1\linewidth]{WK1_Results_files/figure-latex/projectoutline-1}

\begin{ltbluebox}
\emph{This workshop addressed project goals:}\\
* Review and assess existing indicators and evaluate their utility to
support management objectives\\
* Identify and develop priority ecosystem and habitat indicators,
including associated targets and thresholds as appropriate, and identify
specific pathways for incorporation into Council management processes

\end{ltbluebox}

\end{col}

\end{cols}

\end{document}
